% Insert as A.52.  Uses the ordinary ghosts of A.44 and the actual E_3 of A.43.
\begin{statement}{A.52}{Proved}{Ordinary ghosts retaining the full class-two lift and any fixed finite third layer}
Run the recursion of A.44 with $n_j=(K+j)!$, where $K\ge4$,
and put $N=K!$.  Let $(A,B)$ be its ordinary infinite-support
pair and let $s_*$ be its distinguished global class-two lift
from A.40.  Then
\[
 \mathbf o_3(s_*)\longmapsto0
   \quad\text{in }H^2\bigl(G_\Q,
          (\widehat\Z/N\widehat\Z)(3)\bigr)^2.
 \tag{A.52.1}
\]
Equivalently, $s_*$ lifts to the actual arithmetic quotient
\[
 E_{3,N}:=E_3/\bigl(N(\Gamma_3\Pi/\Gamma_4\Pi)\bigr)
                   \longrightarrow E_2.
 \tag{A.52.2}
\]
The kernel in (A.52.2) is $(\widehat\Z/N\widehat\Z)(3)^2$;
the class-two quotient is retained integrally.  In particular
the original recursion, with $K=4$, gives such a lift for
$N=24$.  For every prescribed positive integer $M$, there
exists an ordinary infinite-support pair whose distinguished
class-two lift extends to $E_{3,M}$.
\end{statement}
\begin{proof}
Changing the initial factorial in A.44 does not change the
proved induction.  All $n_j$ remain divisible by $24$;
$n_j\mid n_{j+1}$; every prime dividing $n_{j+1}$ is at most
$K+j+1<(n_j+1)^4$; and the initial equation
$3-2=1$ satisfies every required precision.  The same field,
Hensel, finite-field and convergence arguments therefore give
the stated ordinary pair, with infinite distinct supports.
Every new prime $q_j,r_j$ is $1$ modulo $n_{j+1}$, hence
$1$ modulo $24$.

In additive Kummer notation the explicit products give
\[
 A=\kappa(3)+N A_1,\qquad
 B=\kappa(-2)+N B_1,
 \tag{A.52.3}
\]
where $A_1,B_1$ are the completed classes of
$\prod_jq_j^{n_j/N}$ and $\prod_jr_j^{n_j/N}$.
These products converge in the completed multiplicative group,
because $n_j/N$ is eventually divisible by every fixed integer.
Their real signs are positive.  Their reductions modulo two
are finite-support rational square classes: $n_j/N$ is even
for all sufficiently large $j$.  This observation will be used
only at that finite coefficient level.

Choose continuous cocycles $a_0,b_0$ for the point pair
$(\kappa(3),\kappa(-2))$ and let $h_0$ represent its actual
point section in $E_2$.  Using positive real roots for $3$
and the compatible roots of $-2$ gives
$a_0(\sigma)=0$, $b_0(\sigma)=-1$ at the chosen complex
conjugation $\sigma$.  Choose cocycles $a_1,b_1$ for
$A_1,B_1$ with $a_1(\sigma)=b_1(\sigma)=0$; their real
Kummer classes are zero, so a global coboundary makes these
equalities exact.  Put
\[
              a=a_0+Na_1,\qquad b=b_0+Nb_1.
\]
Take any continuous primitive $h$ given by A.38 for these
abelian representatives.  Both $h$ and $h_0$ vanish at
$\sigma$: the class-two cocycle identity at $\sigma^2=1$
reads $2h(\sigma)=0$, since $a(\sigma)=0$, and the coefficient
module is torsion-free.

\emph{An integral central twist aligns the finite primitive.}
Fix $\ell\mid N$ and set $r=v_\ell(N)$.
The difference
\[
               c=(h-h_0)\bmod\ell^r
 \tag{A.52.4}
\]
is a one-cocycle with values in $(\Z/\ell^r)(2)$: the two
cup defects agree modulo $N$.  We show that this actual
cocycle is unramified away from $\ell$.

At a new support prime $q_j\ne\ell$, the inertia coordinates
of $(a,b)$ are $(n_j,0)$, and the ordinary relation gives
$h(T)=n_jv$ with $v\in\Z_\ell$, as in (A.46.14).
At $r_j\ne\ell$ the coordinates are $(0,n_j)$ and
$h(T)=-n_ju$ with $u\in\Z_\ell$.
Both values are divisible by $N$.  The point section at $3$
has zero class-two inertia coordinate at these primes, by
(A.46.13) with both its valuation coordinates zero.
At any prime outside these supports and $\{2,3,\ell\}$,
the same equation gives zero inertia for both sections.

At $2$, when $\ell\ne2$, the classes of $3$ and $A_1$ have
zero local $\ell$-adic Kummer class: they are $2$-adic units,
whose unit group is pro-$2$.  Write
$a_0=(\chi-1)u_0$ and $a_1=(\chi-1)u_1$ locally.
The $b$-valuation of both sections is one, so their inertia
coordinates differ by $-Nu_1$.
At $3$, when $\ell\ne3$, the classes of $-2$ and $B_1$
both vanish locally.  For the second assertion, each $r_j$
is $1$ modulo $3$, so each finite product defining $B_1$
lies in $1+3\Z_3$ and has zero $\ell$-adic Kummer class;
continuity gives the same conclusion for its completed limit.
Write $b_0=(\chi-1)v_0$, $b_1=(\chi-1)v_1$.
The $a$-valuation is one for both sections, and their inertia
coordinates differ by $Nv_1$.  In particular there is no
loss of a power of two from dividing a Frobenius value by
$3-1$.

All discarded prime-to-$\ell$ inertia acts trivially on the
pro-$\ell$ class-two group and has zero cocycle restriction,
as in A.46.  The preceding tame values therefore show that
(A.52.4) itself descends to
$G_S=G_{\Q,\{\ell,\infty\}}$.  It is also zero on real
complex conjugation.

For odd $\ell$, A.27 gives $H^2(G_S,\Z_\ell(2))=0$.
The coefficient long exact sequence shows that every class
in $H^1(G_S,(\Z/\ell^r)(2))$ is the reduction of an
integral class.  At two, the Bockstein of (A.52.4) has zero
real restriction, while A.27 proves that
\[
 H^2(G_S,\Z_2(2))\longrightarrow H^2(G_\R,\Z_2(2))
\]
is injective.  Its Bockstein is therefore zero in this case
as well, and it too lifts integrally.  By A.40 these integral
classes are multiples of $\eta$ at two and three, and zero
at all other primes.  A single $k\in\Z/24$ realizes the two
required multiples.  After replacing $h$ by $h+k\eta$ and
changing it by a central coboundary, we can consequently arrange
\[
                              h\equiv h_0\pmod N.
 \tag{A.52.5}
\]
This preserves the full integral abelian representatives.

\emph{The binomial term at two.}
At an odd coefficient prime dividing $N$, (A.52.3) and
(A.52.5) make the two cochains in (A.43.1) congruent to the
point-section cochains modulo $\ell^r$, since $2$ is invertible.
Their obstruction classes therefore vanish after this reduction.

At two set $r=v_2(N)\ge3$.  The precise binomial discrepancy is
\[
 \binom{a_0+Na_1}{2}-\binom{a_0}{2}
                   \equiv-\frac N2a_1\pmod{2^r},
 \tag{A.52.6}
\]
and the same identity holds for $b$.
All other terms of (A.43.1) agree modulo $2^r$ by
(A.52.3) and (A.52.5).  The resulting difference of obstruction
classes is the image of
\[
 \bigl(A_1\bmod2\bigr)\cup\kappa_2(-2),\qquad
 -\bigl(B_1\bmod2\bigr)\cup\kappa_2(3)
 \tag{A.52.7}
\]
under the coefficient map sending $1\in\F_2$ to
$N/2\in\Z/2^r(3)$ in the chosen Tate coordinates.
This map is Galois-equivariant, because its image has order
two; the appearances of different Tate weights modulo two
therefore give this exact coefficient map.

Both finite-coefficient cup classes in (A.52.7) are zero.
We verify every place on their finite-support rational
square-class representatives noted after (A.52.3).
At $2$ these representatives are $1$ modulo $8$, hence
squares in $\Q_2$.  At $3$ they are $1$ modulo $3$, hence
squares in $\Q_3$.  At a prime in either support, that prime
is $1$ modulo $24$, so it splits in $\Q(\mu_{24})$, which
contains both $\sqrt{-2}$ and $\sqrt3$.  Thus the other
factor of the relevant cup is a local square there.
At every remaining odd finite prime both factors are units;
their cup is unramified and zero, since the residue-field
Galois group has cohomological dimension one.  At the real
place the root-class representatives are positive.  The
finite-coefficient Brauer localization injection used in A.25 proves
the global vanishing of (A.52.7).  This argument concerns
finite square classes, so no inverse-limit conclusion is
being drawn from their local vanishing.

Equations (A.52.6)--(A.52.7) prove that the chosen integral
class-two lift has zero weight-three obstruction modulo
$2^r$ too.  Thus its two obstruction classes vanish modulo
$N$ at every coefficient factor.

\emph{Identification with the distinguished lift and the actual
quotient section.}
At $(q,\ell)=(2,3)$ the integral local obstruction group is
$\Z/3$, and at $(3,2)$ it is $\Z/8$ by (A.43.3).
Reduction modulo $N$ is injective on each group, since
$3\mid N$ and $8\mid N$; this follows equally from the
finite-coefficient tame residue complex.  Their integral
local obstructions are therefore zero.  The two valuation
vectors are $(0,1)$ and $(1,0)$, so (A.43.4) forces both
central comparison coordinates in (A.40.4) to be zero.
The chosen integral class-two class is exactly $s_*$.
This proves (A.52.1).

The obstruction for (A.52.2) is the reduction of the actual
defect in A.43.  Its vanishing gives two continuous finite-valued
primitives, which correct the last two coordinates in the
Hall expression (A.43.8) and give a section of (A.52.2).
Finally choose $K$ so that $M\mid K!$.  The map on the
central layer $\widehat\Z/K!\widehat\Z\to
\widehat\Z/M\widehat\Z$ induces the actual quotient
$E_{3,K!}\to E_{3,M}$.  The constructed section descends
through this map.  The underlying ordinary pair still has
infinite support, completing the assertion for every
prescribed $M$.
\end{proof}

The quantifiers here fix the finite modulus before constructing
the ghost.  The statement does not provide a single nonrational
pair lifting through all the quotients (A.52.2).  For each fixed
constructed pair, its full integral class-three obstruction
remains the class calculated in A.46.

A.58--A.59 give a further construction with the stronger
valuation divisibilities needed for one full integral
class-three lift. That construction does not take an inverse
limit of the separately chosen ghosts in the present statement.
